% !TeX program = xelatex
\documentclass[12pt,a4paper]{ctexart}
\usepackage{amsmath,amssymb,bm}
\usepackage{booktabs,array,longtable}
\usepackage{graphicx}
\usepackage[margin=2.5cm]{geometry}
\usepackage{caption}
\captionsetup{font=small,labelsep=quad}
\usepackage{listings}
\usepackage{float}
\lstset{basicstyle=\ttfamily\small,breaklines=true,columns=flexible,frame=single,numbers=left,numberstyle=\tiny}
\usepackage[colorlinks=true,linkcolor=black,citecolor=black,urlcolor=black]{hyperref}
\graphicspath{{../05-图表/}}
\numberwithin{equation}{section}
\renewcommand{\arraystretch}{1.25}

% 中文数字章节编号：一、二、…；二级标题为“章.节”式
\ctexset{
  section = {name = {,、}, number = \chinese{section},
             format = \centering\heiti\zihao{4}\bfseries},
  subsection = {number = \arabic{section}.\arabic{subsection},
                format = \songti\bfseries\zihao{-4}},
  subsubsection = {number = \arabic{section}.\arabic{subsection}.\arabic{subsubsection},
                   format = \songti\bfseries\zihao{-4}},
}

% 图未就位时先占位，就位后自动使用；宽度参数适配图的设计宽度
\newcommand{\figplaceholder}[1]{\fbox{\parbox[c][5.5cm][c]{0.86\linewidth}{\centering 图位：#1}}}
\newcommand{\includefig}[3][0.86\linewidth]{\IfFileExists{../05-图表/#2}{\includegraphics[width=#1]{#2}}{\figplaceholder{#3}}}

% 首页版面：article 的 \@maketitle 无条件输出 {\large\@author\par} 与 1.5em 空行，
% 在 \author{} 为空时白占约 32pt。此处只去掉这一行，使摘要（含关键词）整块落在第 1 页。
\makeatletter
\renewcommand\@maketitle{%
  \newpage
  \null
  \vskip 2em%
  \begin{center}%
  \let\footnote\thanks
    {\LARGE\@title\par}%
    \vskip 1em%
    {\large\@date}%
  \end{center}%
  \par
  \vskip 1.5em}
\makeatother

\title{\textbf{中药材烘干过程的热湿耦合建模、守恒求解与分层验证}\\[2pt]
——基于质量坐标有限体积的四问统一求解}
\author{}
\date{2026 年 9 月 12 日}

\begin{document}
\maketitle

\begin{abstract}
\noindent\textbf{问题与挑战}：中药材烘干以圆柱形药柱为对象（长 25 cm、半径 2 cm，
初温 28 ${}^\circ$C、初始干基含水率 2.55 kg/kg），要求给出预热段与全过程的温度、
水分浓度场，固定几何下的达标时间，以及计入收缩后的达标时间。困难有三处：
三套附录材料参数须按问题分别使用，不得调和或常数化；扩散系数 $D$ 非线性强，
尾段最多缩小 16.8 倍；环境数据仅覆盖 4 h，远短于干燥历时，
半径数据在约 21 h 后进入近平台段。任一处处理不当，达标时间的偏差都会达到数十小时。

\smallskip\noindent\textbf{方法}：四问统一为\textbf{一套热湿耦合方程加四个配置开关}：
由干固体守恒、水分守恒与能量守恒导出，以干物质标签 $q$ 为材料坐标写成质量坐标
形式，收缩几何下形式不变且无伪对流项；求解用守恒型有限体积（表面半单元串联阻力、
后向欧拉加局部外推、Picard--Newton 切换、自适应步长），并以题给参数构造判据族
（$Le$、$Bi_m$、$Bi_h$）与参考时标，用于简化与验证。

\smallskip\noindent\textbf{关键结果}：问题 1 温度场与解析解逐点比较，最大偏差
$5.9\times10^{-7}\,^{\circ}$C；问题 3 达标时间 $\bm{t_f=57.1799}$ \textbf{h}
（Richardson 外推 57.1636 h，常系数参考时标 15.70 h 通过）；问题 4 达标时间
$\bm{t_f=50.5369}$ \textbf{h}（瞬时首模近似参考 19.02 h 通过，几何由内生收缩模型预测，
附件 2 实测半径作交叉验证，偏差 $-0.56\%$）。机制归因：附录 4 与附录 3 的扩散
系数比随含水率由 $1/5.39$ 变化到 $1/2.10$；附录材料参数组合效应使固定半径烘干时间放大
$2.26$ 倍（$+78.57$ h），半径收缩（$R^2$ 降 64\%）则加速 $60.9\%$，两者相抵后问题 4
反而快 $11.6\%$；敏感性分析显示边界汇取值（$C_e$）是唯一能改变结论量级的单点（$+22.6$ h）。

\smallskip\noindent\textbf{验证与创新}：验证分六层，每层都说明能排除与不能排除的错误
类型；解析解与数值解是两种不同数学结构的互证；已检查因素的累计偏移量级约 1.71 h
（数值误差、环境情景、模型结构差异三类分别报告）。
内生收缩模型是三项物理动机机制的半经验闭合，含 5 个标定参数，以 RMSE 0.047 mm
贴合附件 2 全部形态（\textbf{单向标定贴合度}）；耦合主解全程 RMSE \textbf{0.088 mm}
（满 72 h 真实轨迹：0--达标 0.065 / 达标--72 h 0.127 mm）；
力学驱动扩展 RMSE 0.522 mm、$t_f=51.1667$ h，结构性上限与材料缺口均已定量说明。

\smallskip\noindent\textbf{关键词：} 热湿耦合\quad 质量坐标\quad 守恒有限体积\quad 判据族\quad 分层验证\quad 中药材干燥
\end{abstract}

\section{问题重述}

\subsection{问题背景}
中药材干燥是药材加工的关键环节：水分过高会导致霉变，干燥不足或过度都影响品质。
本题药材为圆柱形（长 $L_0=25$ cm、半径 $R_0=2$ cm），初始温度 $28\,^{\circ}$C、
初始干基含水率 2.55 kg/kg，置于烘房中干燥；烘房空气温度与空气水分浓度由附件 1
给出（0--14400 s、60 s 一步），药材失水伴随外半径收缩（附件 2，0--72 h、
1800 s 一步，由 2.0 cm 缩至 1.198 cm 后进入近平台段）。烘干达标定义为药材各处
水分浓度（干基含水率）首次全部低于 0.15 kg/kg。

\subsection{问题提出}
四个问题依次为：
\begin{enumerate}
  \item[\textbf{问题 1}] 预热平衡阶段（0--1800 s）药材内温度与水分浓度的时空分布
        （附录 2 的材料参数），按题给表格格式交付（result1.xlsx）；
  \item[\textbf{问题 2}] 整个烘干过程的温度与水分浓度分布（附录 3 的材料参数，result2.xlsx）；
  \item[\textbf{问题 3}] 固定几何下首次全部达标的烘干时间（表 5、result3.xlsx）；
  \item[\textbf{问题 4}] 计入失水收缩（附件 2 外半径、附录 4 的材料参数）后的达标时间
        （表 6、result4.xlsx）。
\end{enumerate}
总体思路是不建立四个孤立模型，而用\textbf{一套热湿耦合控制方程加四个
配置开关}（材料参数来源、热湿耦合、几何是否演化、输出形式）统一求解：问题 1 是
问题 2 的解耦特例，问题 3 是问题 2 加阈值判据，问题 4 是移动域推广；对收缩几何，
先以物理启发的半经验闭合由内而外地预测半径（附件 2 降为验证数据），再以力学驱动
扩展模型讨论收缩的机制边界。几何模型与边界条件示意见图~\ref{fig:几何示意}。

\begin{figure}[htbp]
\centering
\includefig[\linewidth]{fig17_几何模型与边界条件示意图.pdf}{几何模型与边界条件示意}
\caption{几何模型与边界条件。长 25 cm、半径 2 cm 的圆柱药柱，
初始 $U_0=2.55$ kg/kg、$T_0=301.15$ K；柱面与环境 $T_a(t)$、$C_{\mathrm{env}}(t)$
经对流换热 $h=25$ W/(m$^2\cdot$K) 与表面水分交换 $h_m=8\times10^{-7}$ m/s 耦合；
轴心对称、端面忽略（一维径向假设，见第三章）。}
\label{fig:几何示意}
\end{figure}

\section{问题分析}

\subsection{四问的统一结构与难点}
四问共享同一物理过程：内部导热/导湿、表面对流边界、材料参数强非线性。难点有三：
\textbf{(i)} 三套附录材料参数（附录 2/3/4 的 $\rho,c_p,k,D$ 互不相同）必须按问题
分别使用，不得调和、平均或常数化。例如附录 4 与附录 3 的扩散系数比
$0.175\,\mathrm e^{0.15/U}$ 在高含水端约 $1/5.4$（图~\ref{fig:扩散系数}），
把任一材料参数常数化后结果仍看似合理（仍落在 2--3 天），却违反题目指定的建模方式；
\textbf{(ii)} $D\propto\mathrm e^{-\alpha/C}\mathrm e^{-3850/T}$ 对局部含水率与温度
强非线性，尾段 $D$ 最多缩小 16.8 倍，形成低扩散率"干壳"，决定达标时间的尾段行为；
\textbf{(iii)} 环境数据仅 4 h、半径数据 72 h 而干燥需 50 h 以上，边界外推与
收缩闭合的取值选择直接改变结论量级。

\begin{figure}[htbp]
\centering
\includefig[\linewidth]{fig11_扩散系数与时间尺度.pdf}{扩散系数与特征时间尺度}
\caption{扩散系数非线性与特征时间尺度：(a) 三套附录材料参数的
$D$–$C$ 半对数曲线——附录 2 只依赖 $C$、不含 $T$，附录 3/4 为 $D(C,T)$；附录 4 全程低于附录 3，
尾段（$C\to0.15$）$D$ 缩小 16.8 倍（该倍率按附录 3 材料参数公式算出、图上未单独标注）；
(b) 各参考时标：常系数首模参考 15.7021 h、膜界 11.1720 h、问题 4 瞬时首模时间积分
19.016 h、固定 $R_0$ 版 37.61 h、$R_{\min}$ 常值版 15.95 h——主值 57.18/50.54 h
均远大于各时标。}
\label{fig:扩散系数}
\end{figure}

\subsection{各问分析与建模路线}
\textbf{问题 1（预热段场）}：0--1800 s 含水率变化很小，热湿近似解耦（$D$ 不含 $T$、
按附录 2 为常数材料参数）。可建立线性热传导的解析谱解（圆柱 Robin 特征展开 + 时变环境
Duhamel 卷积）作为独立参考，数值解与之逐点比较，这是排除离散、边界、时间积分
三类实现错误的最佳窗口（图~\ref{fig:对拍}）。

\textbf{问题 2/3（全程场与达标时间）}：必须保留双向耦合（$D$ 含 $C,T$）与空间分布。
判据 $Bi_m\in[1.1805,47.7480]$ 表明内部扩散主导，任何"集总/平均含水率"处理
都会把达标时间低估 46.03 h（$-80\%$），不可接受。路线：质量坐标守恒有限体积 +
自适应时间积分，以解析退化、守恒恒等式、网格收敛多层级验证；达标判据用未舍入值
在重建场上判定（中心为最慢点，逐步数值扫描确认）。

\textbf{问题 4（收缩下达标时间）}：几何演化使方程成为移动域问题。在材料坐标下
仿射收缩 $r=R(t)\sqrt q$ 保持方程同形、无伪对流项；半径本身如何给定是核心建模
决策。本文以"失水收缩 + 毛细-膨压塌陷 + 结皮-成孔停滞"三项物理动机机制的
半经验闭合由湿度场逐步\textbf{内生预测} $R(t)$（不读附件 2），附件 2 实测半径
降为验证数据；并以力学驱动扩展模型（有限应变径向平衡 + 可标定保水本构）
讨论收缩的机制结构与精度边界（第 5.5 节）。

\subsection{输入数据与预处理}
附件 1 的环境时程（241 点）用分段线性插值取值，4 h 后采用末值保持外推
（$T_a=50.165\,^{\circ}$C、$C_{\mathrm{env}}=0.04986$；其他取值方式作对照，纳入敏感性分析）；
附件 2 半径用 PCHIP 保形插值（已核查半径为正、单调不增），72 h 后末值延拓，
数据与插值行为见图~\ref{fig:输入数据}。

\begin{figure}[htbp]
\centering
\includefig[\linewidth]{fig1_输入数据.pdf}{输入数据}
\caption{输入数据：附件 1 烘房温度与空气水分浓度
（0--4 h、60 s 步）及末值保持外推；附件 2 外半径（0--72 h）的三段结构
（急缩/缓缩/平台）与约 21 h 的近平台起点。}
\label{fig:输入数据}
\end{figure}

\section{模型假设}
\begin{enumerate}
  \item \textbf{一维径向、轴对称，忽略端面与轴向}（理由：端面/侧面积比 $R/L=8\%$、
        轴向扩散路径 $L/2=12.5$ cm 是径向的 6.25 倍，扩散时间尺度约慢 $6.25^2\approx39$
        倍；本文未做有限长圆柱对照，端面效应不给定量误差界）。
  \item \textbf{湿边界为声明的有效闭合}：把附件 1 的 $C_{\mathrm{env}}(t)$ 直接当作与
        $h_m$ 配套的 Robin 驱动势 $-D\partial_rU=h_m(U_s-C_{\mathrm{env}})$
        （理由：空气与固体含水率基准不同，严格闭合需吸附等温线或水活度关系，
        题目未给；该闭合是本文最大单点依赖，其不确定性由 $C_e=0.1366$ 敏感性
        $+22.6261$ h 量化，见第七章）。
  \item \textbf{热方程取无潜热的有效显热近似}（理由：这是本文声明的有效基线模型；
        潜热通量与显热通量同量级（通量比 2.06--2.49，以第一问 1800 s 点的
        表面量后验复算比值达 5.54）、非小扰动，其影响以表面相变汇对照情景界定：
        情景下 $t_f$ 由 57.1799 h 变为 62.5694 h（$+5.3889$ h），见 5.3 节；
        更真实的闭合需要题目未给的吸附/解吸数据）。
  \item \textbf{材料参数逐点、按局部瞬时值取值}：$\rho,c_p,k,D$ 按附录公式在局部 $C$、
        局部 $T$（开尔文）取值；$h=25$ W/(m$^2\cdot$K)、$h_m=8\times10^{-7}$ m/s
        自附录 2 延用到问题 2--4（声明的参数延用假设）。
  \item \textbf{环境外推}：$t>14400$ s 取末值保持（对照：线性外推、末 1 h 均值，
        差异 $+1.62$ h 以内，已纳入敏感性分析）。
  \item \textbf{问题 4 几何}：长度不变、内部等比仿射收缩 $r=R(t)\sqrt q$
        （附件 2 只识别外半径，等比是声明的最简闭合）；$R(t)$ 由半经验收缩闭合
        内生预测，72 h 后保持 1.198 cm。
  \item \textbf{$\rho(C)$ 不作质量/体积闭合}：$\rho$ 只用于有效体积热容
        $a_{\mathrm{eff}}=\rho c_p$ 与通量算子 $\rho_{sv}=\rho/(1+C)$
        （理由：附录 3/4 的 $\rho(C)$ 与附件 2 的 $R(t)$ 在干固体质量守恒下
        不相容，要求 $R\ge1.301$ cm 而数据终值为 1.198 cm）。
\end{enumerate}

\section{符号说明}
\begin{table}[H]
\centering
\caption{主要符号（首次进入正文处另有解释；附录 B 给出全部参数取值）。}
\label{tab:符号}
\small
\begin{tabular}{lll}
\toprule
符号 & 含义 & 单位 \\
\midrule
$q$ & 材料坐标（干物质标签），$q=(X/R_0)^2$ & — \\
$U(q,t)$ & 干基含水率（水分浓度） & kg/kg \\
$T(q,t)$ & 温度（开尔文） & K \\
$D(C,T)$ & 扩散系数（附录 2/3/4 公式） & m$^2$/s \\
$\rho,\ c_p,\ k$ & 密度、比热容、导热系数（逐点取值） & kg/m$^3$、J/(kg$\cdot$K)、W/(m$\cdot$K) \\
$a_{\mathrm{eff}}$ & 有效体积热容 $\rho c_p$ & J/(m$^3\cdot$K) \\
$h,\ h_m$ & 表面对流换热系数、表面传质系数 & W/(m$^2\cdot$K)、m/s \\
$T_a(t),\ C_{\mathrm{env}}(t)$ & 环境温度、环境水分浓度（附件 1） & K、kg/kg \\
$R(t)$ & 药柱外半径（问题 4 为模型输出） & m \\
$t_f$ & 达标时间（$\max_q U<0.15$ 首次成立） & h \\
$Le,\ Bi_m,\ Bi_h$ & 判据族：$\alpha/D$、$h_mR/D$、$hR/k$ & — \\
$J$ & 体积变形梯度（力学扩展） & — \\
$S$ & 液相饱和度（力学扩展） & — \\
$p_c,\ \bar p_c$ & 毛细压及其储能平均（力学扩展） & （模型单位） \\
\bottomrule
\end{tabular}
\end{table}

\section{模型建立与求解}

\subsection{统一控制方程与质量坐标}
\textbf{守恒律与材料坐标}：从干固体守恒与水分守恒出发（推导见附录 A），
以干物质标签 $q\in[0,1]$（$q=(X/R_0)^2$，$X$ 为初始物质点）描述径向一维
热湿耦合过程（毛细多孔介质热湿传递的经典框架见~\cite{luikov}）；在干物质密度空间均匀与仿射收缩 $r=R(t)\sqrt q$ 下，
方程与固定域同形、无伪对流项：
\begin{equation}
U_t=\partial_q\!\left(\frac{4qD}{R(t)^2}U_q\right),\qquad
a_{\mathrm{eff}}\,T_t=\partial_q\!\left(\frac{4qk}{R(t)^2}T_q\right).
\label{eq:质量坐标}
\end{equation}
边界（$q=1$，移动表面 Robin 条件形式不变——跳条件中几何运动项恒等相消）：
\begin{equation}
-D\,\partial_rU=h_m\bigl(U_s-C_{\mathrm{env}}(t)\bigr),\qquad
-k\,\partial_rT=h\bigl(T_s-T_a(t)\bigr),
\label{eq:边界}
\end{equation}
$q=0$ 为正则对称；初值 $U\equiv2.55$、$T\equiv301.15$ K。
式~\eqref{eq:质量坐标} 与边界式~\eqref{eq:边界} 的全部符号见表~\ref{tab:符号}；
四问由四个开关配置（表~\ref{tab:开关}）：

\begin{table}[htbp]
\centering
\caption{四问 = 一套耦合方程 + 四个开关。}
\label{tab:开关}
\small
\begin{tabular}{lp{0.17\linewidth}p{0.16\linewidth}p{0.17\linewidth}p{0.22\linewidth}}
\toprule
开关 & 问题 1 & 问题 2 & 问题 3 & 问题 4 \\
\midrule
材料参数 & 附录 2（常数，$D=f(C)$） & 附录 3 & 附录 3 & 附录 4 \\
热湿耦合 & 解耦（$D$ 不含 $T$） & 双向耦合 & 双向耦合 & 双向耦合 \\
几何 & 固定 $R=0.02$ m & 固定 & 固定 & $R(t)$ 内生（半经验闭合，附件 2 交叉验证） \\
环境 & 附件 1（0--1800 s） & 附件 1 + 外推 & 同左 & 同左 \\
输出 & 场（表 1/2，result1） & 场（表 3/4，result2） & 决策量 $t_f$（表 5，result3） & 决策量 $t_f$（表 6，result4） \\
\bottomrule
\end{tabular}
\end{table}

\textbf{数值方法}：固定 $q$ 网格单元中心有限体积（N=160 主网格），内面通量共享；
表面用\textbf{半单元内阻与外对流阻串联}，表面值 $U_s,T_s$ 由
$8D_{\mathrm{eff}}(U_s-U_N)/(R^2\Delta q)=2h_m(C_{\mathrm{env}}-U_s)/R$ 联立重建
（$D_{\mathrm{eff}}$ 取半单元中点值，避免强变的 $D\propto\mathrm e^{-0.45/C}$
在端点取值产生"内通量被强制为零"的伪根），不把末单元平均值当作表面值。
时间用后向欧拉加局部外推（一步对两半步，接受 $2\times$半对$-$整步，二阶），
外推态的守恒通量由各后向欧拉恒等式线性组合给出，仍严格守恒；
自适应步长（rtol $=10^{-7}$，误差比控制，失败缩步不截断负值），
尾段步长约 10--15 s；输出 1 s/60 s 网格由步记录线性插值重构
（插值误差折合约 $2.6\times10^{-4}$ h，见第七章）。
非线性用外场 Gauss--Seidel（热场在冻结系数下线性求解，湿场在冻结温度下
Newton/Picard），湿场 Picard 起步、$\max_qU<0.6$ 后切换含 $\partial D/\partial C$
的 Newton（全程 Picard 对照 $\Delta t_f=1.7\times10^{-10}$ h，可忽略）。
全部交付由 \texttt{python run\_all.py} 一条命令产出，无人工干预；
连续两次运行 result1--4 的 md5 完全一致（附录 D 列出各文件的哈希值）。

\subsection{问题 1：预热平衡段的场（0--1800 s）}
\textbf{模型}：附录 2 为常数材料参数、热湿解耦，热方程为线性圆柱 Robin 问题；
湿分在 1800 s 内变化小但非零，仍按同一框架求解。
\textbf{求解与局部检验}：数值解与解析谱解（圆柱 Robin 特征展开 + 时变环境
Duhamel 卷积，稳态减法消除级数尾巴，解析参考自身误差 $<10^{-8}\,^{\circ}$C）
逐点比较：单元中心最大偏差 $5.9\times10^{-7}\,^{\circ}$C（N=320），
收敛阶 2.00/1.99/1.99（图~\ref{fig:对拍}）——据此排除离散、边界处理、时间积分
三类实现错误（材料参数取值错误不在此层排除之列）。
\textbf{结果}：温度与水分浓度径向分布见图~\ref{fig:温湿剖面}（色阶按时间递增）、
时空云图见图~\ref{fig:温度云图}；特征位置的时程与偏差、
表面通量与表面扩散系数见图~\ref{fig:表面通量}；交付表见表~\ref{tab:表1}、
表~\ref{tab:表2}（result1.xlsx）。

\begin{figure}[htbp]
\centering
\includefig[\linewidth]{fig12_温度场时空云图.pdf}{温度场时空云图}
\caption{温度场时空演化：时间 × 径向位置的温度云图显示，
表面在约 600 s 内贴近环境、中心滞后明显，1800 s 末全场温差收敛到
约 $1\,^{\circ}$C 量级；等温线的疏密即热透进程。}
\label{fig:温度云图}
\end{figure}

\begin{figure}[htbp]
\centering
\includefig[\linewidth]{fig14_表面通量与表面扩散系数.pdf}{表面通量与表面扩散系数}
\caption{表面通量与表面扩散系数。(a) 表面水分
通量 $j_s$ 与表面扩散系数 $D_s=D(C_s)$ 的时程：通量初期最大而后衰减，
$D_s$ 随表面变干持续下降；(b) $D_s$ 与 $C_s$ 的对应关系，本过程的 $C_s$
只覆盖高含水端，故曲线平缓；$D\propto\mathrm e^{-0.89/C}$ 在全量程上的
强非线性见图~\ref{fig:扩散系数}(a)。}
\label{fig:表面通量}
\end{figure}

\subsection{问题 2/3：全程场与固定几何达标时间}
\textbf{模型}：附录 3 的材料参数、双向耦合、固定几何；达标判据
\begin{equation}
t_f=\inf\{t:\ \max_{q\in[0,1]}U(q,t)<0.15\ \text{kg/kg}\},
\label{eq:达标}
\end{equation}
用未舍入值在全域重建场上判定（判据式~\eqref{eq:达标}；中心是否最慢点不作公理，逐步数值扫描确认：
全部时刻最大值均位于中心列，容差 $10^{-12}$）。
\textbf{结果（场）}：温湿剖面随时间的演化（色阶按时间递增，viridis）见图~\ref{fig:温湿剖面}，
水分场时空云图与干壳/湿核结构见图~\ref{fig:水分云图}；干燥锋面自表面向内推进、
表面先干、中心最慢。随表面变干，$D\propto\mathrm e^{-0.45/C}$ 最多缩小 16.8 倍，
形成低扩散率"干壳"，中心残余水分只能穿过它排出——尾段 $\ln(U_{\max}-C_e)$
呈指数衰减（$R^2=0.9972$），衰减率对应剖面加权等效
$D\approx2.5\times10^{-10}$ m$^2$/s（后验尾段诊断，见第六章）。

\begin{figure}[htbp]
\centering
\includefig[\linewidth]{fig3_温湿剖面.pdf}{温湿剖面}
\caption{温湿剖面演化。代表时刻的径向温度与水分浓度剖面，
色阶按时间递增（viridis）：表面先干、中心最慢、尾段拖长。}
\label{fig:温湿剖面}
\end{figure}

\begin{figure}[htbp]
\centering
\includefig[\linewidth]{fig13_水分场时空云图.pdf}{水分场时空云图}
\caption{水分场时空演化：时间 × 径向位置的水分浓度
云图显示，干燥锋面内移、表面快速进入低 $D$ 区形成干壳，中心湿核维持到 40 h 以后；
等值线为 0.15 kg/kg 达标线，色阶取对数以便分辨干壳与湿核，中心最后穿越。}
\label{fig:水分云图}
\end{figure}

\textbf{结果（达标时间）}：数值上报告\textbf{穿越区间加右侧安全端点}：
问题 3 为 $[205837.7$, $205847.7]$ s，$\bm{t_f=57.1799}$ \textbf{h}
（插值 57.1795 h，常系数参考时标 15.70 h 通过）。
收敛序列（N=40/80/160）为 57.249/57.201/57.180 h，
Richardson 外推 57.1636 h，主值偏差 0.016 h 即空间误差量级。
四位小数舍入后的 60 s 网格判定为 57.2333 h，与未舍入值差 0.053 h
（另一种数据使用方式的差异，见第七章）；阈值 0.15 干基 $=13.04\%$ 湿基，
与药材水分限度的通行标准（约 13\%，未核对具体药典条文与版本卷次）相符~\cite{pharmacopoeia}。
交付表见表~\ref{tab:表3}、表~\ref{tab:表4}（result2.xlsx）与
表~\ref{tab:表5}（result3.xlsx）；达标穿越见图~\ref{fig:穿越}。

\textbf{潜热对照情景（界定无潜热基线的适用性）}：第三章假设 3 的无潜热闭合
用情景检验界定：表面相变汇 $-k\partial_rT|_R=h(T_s-T_a)+L_vj_w$，
$j_w=\rho_{d,s}h_m(U_s-U_e)$（$\rho_{d,s}=(650+128U_s)/(1+U_s)$ 为表面局部
干物质密度，$L_v=2.4\times10^6$ J/kg 近似值），同一套附录 3 材料参数、同一边界数据、
同一数值框架——$t_f$ 由 57.1799 h 变为 \textbf{62.5694 h}
（$\bm{+5.3889}$ \textbf{h}，$+9.4\%$），全程最大逐时表面/中心温差
$29.6/29.2\,^{\circ}$C、最大平均含水率差 $0.322$ kg/kg。故无潜热基线的 $t_f$ 偏差量级约 5.4 h；
该情景自身闭合同样不完备（情景中表面温度一度降至 $13.4\,^{\circ}$C、
低于湿球温度，暴露出 $h_m$ 与能量边界的不协调），情景值不进入主交付结果。

\subsection{问题 4：收缩移动域与内生收缩模型}
\textbf{几何与移动域}：仿射收缩 $r=R(t)\sqrt q$ 下方程同形（式~\eqref{eq:质量坐标}），
$R(t)$ 为待确定的几何函数。\textbf{附件 2 的三段结构}（A 急缩 0--3.5 h 占 61\%、
B 缓缩 3.5--21 h、C 近平台段约 21 h 起）不是单一失水收缩曲线：
单调映射总能构造出来，但理想等容收缩这条特定基线
$R_0\sqrt{(1+\bar U)/(1+C_0)}$ 不能解释全部数据——其反演平均含水率
（0.278 kg/kg）比模型同刻值（0.149 kg/kg）更低（基线反演的等效含水率更低），
并存在成孔缺口 $0.129\to0.155$ kg/kg 与数据约 16 h 交叉的滞后现象
（附件 2 覆盖 0--72 h，$t_f$ 落在 72 h 观测区间内），见图~\ref{fig:收缩机制}。

\begin{figure}[htbp]
\centering
\includefig[\linewidth]{fig6_收缩机制分析.pdf}{收缩机制分析}
\caption{附件 2 的收缩机制分析。附件 2 呈三段复合结构
（A 急缩 61\%/B 缓缩/C 平台），与约 16 h 滞后交叉；零拟合的"理想收缩+玻璃化冻结"
半经验闭合已能结构性给出减速与平台（均值触发那一条：平台 1.169 cm、偏差 $-2.4\%$、
冻结时刻偏差 $+12\%$），
早期塌陷与后期成孔对应两组机制。数据源：问题 4 主解湿度场与附件 2 实测半径。}
\label{fig:收缩机制}
\end{figure}

\textbf{半经验收缩闭合（问题 4 主模型）}：
\textbf{问题 4 主模型 = 物理启发的半经验收缩闭合内生预测半径}：三项物理动机机制叠加，
一项一机制、总参数 5 个：
\begin{equation}
R(t)=\underbrace{R_0\sqrt{\frac{1+\bar U(t)}{1+C_0}}}_{\text{失水收缩}}\cdot
\underbrace{\bigl[1-\varepsilon_c\,g_c(\bar U)\bigr]}_{\text{毛细-膨压塌陷}}\cdot
\underbrace{\bigl[1+\varepsilon_p\,g_p(C_s)\bigr]}_{\text{结皮-成孔停滞}}
\label{eq:收缩闭合}
\end{equation}
(i) 失水收缩（零参数基线）：水分守恒 + 固体/水等密度声明；
(ii) 毛细-膨压塌陷：$g_c(x)=(1-x)^{a}x^{b}$（$x=(C_0-\bar U)/C_0$，峰值归一），
应变量级有毛细压实估计支撑（bracket $[0.028,0.580]$ cm 覆盖观测超额
0.182 cm @2 h、半径应变 9.1\%）；
(iii) 结皮-成孔停滞：表面 $C_s$ 穿越 $C_{\mathrm{glass}}\approx0.21$
（Gordon--Taylor 结构典型参数情景值~\cite{saravacos}，列入参数不确定清单）后收缩停滞、
继续失水转化为孔隙，
$g_p=\mathrm{clip}\bigl((C_{\mathrm{glass}}-C_s)/(C_{\mathrm{glass}}-C_{s,\min}),0,1\bigr)^{k_p}$，
$C_{s,\min}=0.0488$ kg/kg 为完整轨迹读得的归一化常数（携带轨迹信息，属标定的一部分）。
$g_c,g_p$ 是带物理动机的经验形状函数，并非由应力平衡推出，
本闭合是\textbf{半经验}结果；5 个参数由最小二乘对附件 2 标定
（不是由任何公式"算出"），2 个量（$C_{\mathrm{glass}}$、$\rho_{sk}=1468$）
取情景值/典型范围值不参与拟合（参数表见附录 B，标定细节见附录 C）。
该结构\textbf{不自动保证半径进入平台}；早期超额收缩的正确分解是
孔隙塌陷（$\le9.94\%$，已用尽）+ 基体黏弹压密的组合——后者在零初始孔隙
归一化与组分体积可加假设下对应约 31\% 固体应变，但该比例依赖初始质量、
真实固体密度与组分体积假设，属条件性诊断而非独立测量；"负孔隙"是
"骨架密度 1468 + 零初始孔隙基线"后处理的\textbf{伪影}，不是压密证据
（按体积审计，初始孔隙率峰值需求为 17.26\% @ 2 h；把题给 $\rho(2.55)=989.5$
当堆密度时可用值仅 9.94\%、缺口 7.33 个百分点）。

\textbf{标定与验证}：拟合优度 RMSE $=0.047$ mm（附件 2 量化分辨率 0.01 mm 的
约 5 倍）、最大偏差 0.162 mm，滞后交叉（预测 17.0 h、即约 17 h；数据约 16 h）、
平台（1.1969 vs 1.198 cm）、三段形态全部复现（图~\ref{fig:内生贴合}）；
孔隙份额转正于 9.5 h 后单调增加（成孔），末期孔隙率 30.9\%。
过拟合自检表明：标定依赖链
$R_{\mathrm{obs}}\to(\bar U,C_s)_{\text{主解}}\to\text{标定}\to R_{\mathrm{pred}}$
存在循环（湿度场由附件 2 实测几何生成）；奇偶折半留出法
（训练/留出 0.045--0.050 mm）只检验同一光滑曲线上的\textbf{插值稳健}；
真正的独立验证需要\textbf{完整耦合标定}（每组候选参数从相同初值重解），
本文当前是"借助实测几何生成湿度场再标定"的近似。

\begin{figure}[htbp]
\centering
\includefig[\linewidth]{fig7_内生收缩贴合.pdf}{内生收缩贴合}
\caption{内生收缩模型对附件 2 的贴合与验证：
三项物理动机机制的半经验闭合（式~\eqref{eq:收缩闭合}，5 个标定参数 + 2 个
非拟合量）复现附件 2 的全程形态：RMSE 0.047 mm（量化分辨率的 5 倍）、滞后交叉
（17.0 h vs 约 16 h）、平台（1.1969 vs 1.198 cm）、三段形态全对；
残差子图显示偏差有界在 $\pm0.16$ mm（单向拟合指标）。}
\label{fig:内生贴合}
\end{figure}

\textbf{耦合主解（附件 2 降为验证数据）}：主结果采用完全内生的收缩——
$R$ 由湿度场逐步预测（每步以当前 $\bar U,C_s$ 代入式~\eqref{eq:收缩闭合}），
不读附件 2，得 $\bm{t_f=50.5369}$ \textbf{h}（result4 内生版、3033 行；未舍入右端点，
穿越区间 $[181917.7,181932.7]$ s；瞬时首模近似参考 19.02 h 通过）；
附件 2 外生几何版 $t_f=50.8230$ h 作交叉验证（偏差 $-0.286$ h、$-0.56\%$，
属模型间对照，不构成独立实验验证）。两个精度指标须分清：
\textbf{0.047 mm 是外生湿度轨迹驱动的单向拟合指标}；
\textbf{耦合主解自身指标}为满 72 h 真实轨迹的全程 RMSE \textbf{0.088 mm}
（0--达标 \textbf{0.065} / 达标--72 h \textbf{0.127} mm；
主解积分满 72 h、首次达标单独记录，无常值填充）；
末期 $R_{\mathrm{pred}}(72\,\mathrm h)=1.1836$ vs 附件 2 的 1.198 cm，
60--72 h 再降 0.031 mm（末端近似平台但有残余缓降，与结构限定一致）。
半径预测所用的代数闭合残差（在 $R_{\mathrm{ideal}}$ 下求 $C_s$ 再代入修正式的
未收敛残差）max 0.003 mm；末单元温度与重建表面温度之差 max 0.018 ${}^\circ$C，
对 $t_f$ 的影响远在 0.01 h 以下。
收敛序列 50.5450/50.5370/50.5369 h 的相邻差（28.8 s、0.41 s）后者已小于
尾段约 15 s 的步宽，表观收敛阶不足以说明问题；改用连续事件定位的 $t^*$ 序列
50.5422/50.5355/50.5362 h（N 由 80 增至 160 时变化 0.0007 h，作为空间误差量级，见第七章）。
中心/表面演化与半径收缩见图~\ref{fig:中心表面}，截面演化见图~\ref{fig:截面}；
交付表见表~\ref{tab:表6}（result4.xlsx）。

\begin{figure}[htbp]
\centering
\includefig[0.86\linewidth]{fig16_中心表面演化与半径收缩.pdf}{中心/表面演化与半径收缩}
\caption{中心/表面含水率演化与半径收缩。(a) 中心与表面
含水率曲线 + 0.15 kg/kg 阈值虚线 + $t_f=50.5369$ h 达标竖线：表面先到干态、
中心最后达标；(b) 内生半径 $R(t)$ 在 50.5 h 进入近平台，60--72 h 仅再降 0.031 mm
（模型已积分至 72 h，曲线止于 $t_f=50.54$ h）。}
\label{fig:中心表面}
\end{figure}

\begin{figure}[htbp]
\centering
\includefig[\linewidth]{fig9_截面演化.pdf}{截面演化}
\caption{问题 4 的截面演化。6 个代表时刻的圆盘截面
热图（色标为水分浓度，取对数色阶以便分辨低含水端的壳核差异，
色标上 0.15 一刻即题目达标阈值）显示，干燥自表面向中心推进、外半径同步收缩并进入平台：
附件 2 实测的平台半径为 1.198 cm，模型内生预测的末端半径为 1.19 cm
（各面板标题给出该时刻的模型半径）；图底的时间轴带给出内生半径 $R(t)$
与六个采样时刻在曲线上的位置——六个盘的时间间隔是 5/10/10/15/10 h，
而半径降幅集中在前两盘：$R$ 在头 5 h 由 2.00 降到 1.42 cm（占总降幅 0.81 cm 的 72\%），
20 h 后即进入平台，故"等间隔取样、等宽排版"下视觉变化前快后慢；
中心为最慢点，$t_f=50.5369$ h 时中心首次低于 0.15 kg/kg。}
\label{fig:截面}
\end{figure}

\textbf{与问题 3 的大小关系}（三组对照，机制归因必须定量）：
\begin{itemize}
  \item \textbf{附录材料参数组合效应}：附录 4 与附录 3 的 $D$ 之比
        $0.175\,\mathrm e^{0.15/U}$ 由 $1/5.39$（$U=2.55$）变到 $1/2.10$（$U=0.15$）；
        固定半径下 $t_f$ 从 57.18 h 放大到 129.10 h（$\times2.26$，$+78.57$ h）——
        该对照同时改变了全部热物理参数，不能严格归因于单个扩散系数倍率。
  \item \textbf{收缩效应}：$R^2$ 降 64\%（扩散时间尺度缩至 0.36）、$2/R$ 换热换质
        面积比增 67\%——把 129.10 h 拉回 50.54 h（$-60.9\%$），
        收缩是本问题最大的单点加速因素。
  \item \textbf{净效应}：材料参数组合效应带来的减速（$\times2.26$）被收缩加速反超，
        问题 4 反而比问题 3 快 $11.6\%$（$+6.64$ h 对固定半径对照）。
\end{itemize}

\subsection{力学驱动扩展：收缩的力学机制}
本节把半经验闭合升级为"由守恒律和有限应变力学约束、可标定材料本构的内生收缩模型"，
作为\textbf{扩展研究}：主线（5.4 节半经验闭合）是当前对附件 2 精度最高的可用模型；
本节的力学模型是严谨骨架，用途是说明收缩的机制结构与适用边界，并定量指出
走向物理预测还缺哪些数据。缺保水与力学实测数据，结论限于机制结构与情景分析。

\textbf{框架}：材料点 $X$ 的现时位置 $r(X,t)$ 为未知量，有限应变非仿射几何
$\lambda_r=r_X$、$\lambda_\theta=r/X$、$\lambda_z=1$、$J=\lambda_r\lambda_\theta$；
孔隙非负直接进入方程：$g=J-s_0-w\ge0$ 由乘子 $\zeta$ 以互补条件活动集实施；
毛细储能 $\psi_{\mathrm{cap}}=(J-s_0)\int_S^1p_c\,ds$，保水本构取幂律
$p_c=p^*(1-S)/S^\beta$（$\bar p_c$ 有解析式，与数值积分一致到 $5\times10^{-16}$）；
单松弛对数主应变黏弹骨架（内变量 $Z$，后向欧拉局部消元，每个时间步只提交一次）；
初始平衡 $e_*=\bar p_c(S_0)/K_\infty$ 使初始总应力为零；
总应力 $\sigma_i=H_i/J+\bar p_c-\zeta$，准静态径向平衡弱式求解、表面 $\sigma_{rr}=0$
——\textbf{外半径 $R(t)$ 是输出而非输入}；与传质以 $U_t=\partial_q(D/r_q^2U_q)$
衔接，每步按力学、几何系数、传质的顺序迭代耦合；热算子按
$a_{\mathrm{eff}}T_t=(1/(rr_q))\partial_q(rk/r_qT_q)$ 的有限体积形式离散。
判别实验证明这是\textbf{真正的全局应力传递}而非逐点体积平衡（只干燥最外层时，
湿内部 $J$ 保持 1.000 但径向应力全域均匀传递为 $-0.047$、表皮环向应力 $+4.188$、
表面节点内移 0.042 cm）。

\textbf{正确性验证}（退化检查，表~\ref{tab:力学退化}）：孔隙非负全程成立、
体积衡等式达机器精度、热算子与参考求解器的比较结果一致（常数材料参数、无失水、阶跃环境下
同一后向欧拉步：N=16/80/160 最大场差 $0.0084/0.0023/0.0013\,^{\circ}$C，
随加密单调下降）、牛顿步全部收敛（最大平衡残差 $9.5\times10^{-11}$）。

\begin{table}[htbp]
\centering
\caption{力学求解器的退化/正确性检查（判定列全为通过）。}
\label{tab:力学退化}
\small
\begin{tabular}{p{0.26\linewidth}p{0.26\linewidth}p{0.31\linewidth}l}
\toprule
检查 & 应出现的结果 & 实测 & 判定 \\
\midrule
不干燥 + 初始总应力平衡 & 尺寸无人为跳变 & $\max|\Delta r|=0.000$ & $\checkmark$ \\
关闭黏弹分支 & 恢复纯弹性问题 & 运行正常（$R_{\mathrm{end}}=1.717$ cm） & $\checkmark$ \\
$\tau\ll$ 过程尺度 & 黏弹分支充分松弛 & 与纯弹性差 $5.3\times10^{-15}$ cm & $\checkmark$ \\
$\tau\gg$ 过程尺度 & 内变量近冻结 & $\max|Z|/\max|E|=4.3\times10^{-11}$ & $\checkmark$ \\
指定仿射映射 $r=R\sqrt q$ & 传质系数恢复 $4qD/R^2$ & 相对误差 $4.9\times10^{-15}$ & $\checkmark$ \\
固定外形持续失水 & $g$ 按体积恒等式增长 & 相对误差 $2.2\times10^{-16}$ & $\checkmark$ \\
全程孔隙非负 & 任一时刻任一点 $g\ge0$ & $g_{\min}=+0.1446$（无任何负孔隙） & $\checkmark$ \\
全局体积核验式 & $R^2/R_0^2=s_0+\frac{\rho_{d0}}{\rho_w}\bar U+\int g\,dq$ & 残差 $\le2.7\times10^{-16}$ & $\checkmark$ \\
热算子退化对拍 & 与参考求解器同 BE 步一致，随加密下降 & $0.0084/0.0023/0.0013\,^{\circ}$C & $\checkmark$ \\
\bottomrule
\end{tabular}
\end{table}

\textbf{标定与复现程度}：幂律保水—骨架结构标定到附件 2
（$\beta$ 下界 0.02，4 个起点收敛同一解）：$p^*=1.0\times10^{4}$（触上界，
大 $p^*$ 下轨迹由无量纲形状 $\bar p_c(S)/\bar p_c(S_0)$ 决定，$p^*$ 处于
渐近不变区、不可辨识属预期）、$\beta=0.1059$、$\Pi_\tau=1.0\times10^{-3}$（触下界）。
生产档（满 72 h 积分）\textbf{RMSE $=0.522$ mm}（急缩段 1.193 / 缓缩段 0.875 /
平台段 0.163 mm；分窗 0--达标 0.604 / 达标--72 h 0.212 mm），
末期平台 1.1747 vs 附件 2 的 1.198 cm，60--72 h 仅再降 0.04 mm（末端真平台），
$\bm{t_f=51.1667}$ \textbf{h}（较半经验主线 50.5369 h 偏差 $+1.25\%$）。
奇偶折半留出 0.587/0.591 mm（无过拟合）；块 bootstrap 给出 $\beta$ 的 95\%
置信区间 $[0.095,\,0.91]$；孔隙非负全程成立（$g_{\min}=0.1446$）。
平台由\textbf{吸力饱和平衡}产生（标定参数下表面 Deborah 数
De $\approx1.0\times10^{-3}\ll1$，非黏弹冻结），见图~\ref{fig:力学升级}。

\begin{figure}[htbp]
\centering
\includefig[\linewidth]{fig10_力学升级.pdf}{力学升级}
\caption{力学驱动内生收缩模型：修复后求解器 +
幂律保水重标定（$\beta=0.106$）对附件 2 的拟合结果为 RMSE 0.522 mm
（残差三段：急缩 1.193 / 缓缩 0.875 / 平台 0.163 mm）、$t_f=51.1667$ h；
末端 60--72 h 仅再降 0.04 mm（真平台），平台机制为吸力饱和平衡
（De $\approx10^{-3}\ll1$）。(c) 面板给出该平台机制：归一化平均吸力 $\bar p_c/p^*$
随饱和度 $S$ 下降而单调上升，并在 $S\to0$ 时趋于有限上限
$p^*/[(1-\beta)(2-\beta)]$——末端收缩进入平台即来自这一吸力饱和平衡；
气相体积份额全程非负（$g_{\min}=0.1446$，无任何负孔隙）见表~\ref{tab:力学退化}。}
\label{fig:力学升级}
\end{figure}

\textbf{结构性上限与材料缺口}：在所选幂律保水—简化储能结构内，
$\bar p_c$ 有有限上限 $p^*/[(1-\beta)(2-\beta)]$（$\beta=0.15$ 时仅 $0.636\,p^*$），
且吸力是局部的——只作用于干燥表皮，湿内部以体积刚度抵抗经全局传递来的压缩。
后果：(i) 附件 2 前 21 h 的早期全局塌陷不能由吸力驱动的局部收缩单独给出；
(ii) 末端平衡 $J_{\mathrm{end}}(\beta)$ 由 $J\cdot g(S(J);\beta)=g(S_0;\beta)$ 决定：
$\beta=0.02/0.05/0.10/0.15/0.85$ 时 $J_{\mathrm{end}}=0.368/0.356/0.336/0.319/0.218$
——要到附件 2 的平台 $J=0.36$ 必须 $\beta\lesssim0.05$（但缓缩段过慢），
要让中段够快必须 $\beta\gtrsim0.10$（但末端必过缩），该结构内单一 $\beta$
给不出"急缩快+平台停"；标定的折中 $\beta=0.106$ 末端 $J=0.3343$（vs 0.36）。
此上限是所选本构与储能结构的性质，不排除其他保水曲线或体积驱动机制达到
更高精度。材料缺口：真实保水曲线 $p_c(S,T)$（最大杠杆）、膨压类体积驱动、
干态模量增长/结皮硬化、骨架模量与松弛谱、初始湿质量与真实密度
（体积审计只给 $\rho_{d0}\le256$ kg/m$^3$ 的上限）。
\textbf{结论}：当前精度仍不如半经验主线（0.047/0.088 mm），
在真实保水曲线与力学实测数据到位之前，
不能声称"力学模型复现附件 2"；力学驱动升级是后续正统方向，
届时半经验收缩闭合可作为其低维验证基准。

\section{模型检验}

\subsection{判据族（Le、Bi\_m、Bi\_h）}
沿矩形扫描域 $C\in[0.15,2.55]$、$T\in\{301.15,\,323.315\}$ K、$R=R_0$ 计算
（$\alpha=k/(\rho c_p)$，逐点取附录公式；扫描由 \texttt{run\_m4.criteria\_table()}
执行，区间即程序输出），结果见表~\ref{tab:判据族}。干燥末段表面 $U<0.15$ 的局部状态超出该扫描域，
不在区间覆盖之内。

\begin{table}[htbp]
\centering
\caption{判据族表：定义、扫描区间、阈值语义、简化决策与反验偏差
（反验偏差 = 采用该简化后 $t_f$ 与全耦合解 57.2009 h 的实测差）。}
\label{tab:判据族}
\small
\begin{tabular}{llp{0.16\linewidth}p{0.16\linewidth}p{0.23\linewidth}l}
\toprule
判据 & 定义 & 附录 3 区间 & 附录 4 区间 & 简化决策 & 反验偏差 \\
\midrule
$Le$ & $\alpha/D$ & $10.6860$--$640.6171$ & $31.2068$--$556.1413$ & $Le\gg1$：热场准稳态 $T\approx T_a$ & $0.303$ h \\
$Bi_m$ & $h_mR/D$ & $1.1805$--$47.7480$ & $6.3606$--$100.3742$ & $Bi_m\to0$：集总（内扩散无阻力） & $-46.03$ h \\
$Bi_h$ & $hR/k$ & $1.0353$--$1.9263$ & $1.8964$--$3.4226$ & $Bi_h\to0$：忽略表面热滞后 & $0.075$ h \\
\bottomrule
\end{tabular}
\end{table}

反验结果（每条给出采用该简化后 $t_f$ 的偏移）：
\begin{itemize}
  \item \textbf{$Le\gg1$ 支持热场准稳态}：温度场取 $T_a(t)$ 只求湿分，
        $t_f$ 偏差 $0.303$ h（0.5\%）；本文主模型仍保留双向耦合，此简化只作
        判据反验，且只在无潜热基线内评价，不外推到有蒸发冷却的情形
        （5.3 节潜热情景表面温差达 $29.6\,^{\circ}$C）。
  \item \textbf{$Bi_m$ 区间 $\gtrsim O(1)$ 表明内部扩散主导}：集总模型给出
        $t_f=11.17$ h，比全耦合解快 $46.03$ h（$-80\%$）；可见任何"集总/平均
        含水率"处理都不可接受，本题必须保留空间分布。
  \item \textbf{$Bi_h\sim O(1)$：对本模型达标时间的敏感性小}：令 $h\to\infty$
        实测偏差 $0.075$ h（0.13\%），本文未采用该简化；$Bi_h\sim1$ 本身
        并不表示空间温差小，该结论同样不外推到有蒸发冷却的情形。
\end{itemize}

\subsection{参考时标（常系数与瞬时首模近似）}
取 $D$ 的路径点态最大 $D_{\max}$（$D$ 对 $C,T$ 均单调增），构造常系数 Robin
参考问题，其首模形式为
\begin{equation}
t_{\mathrm{ref}}=\frac{R^2}{D_{\max}\lambda_1^2}\,
\ln\!\frac{A_1\,(C_0-C_e)}{C_{\mathrm{th}}-C_e},
\label{eq:时标}
\end{equation}
其中 $\lambda_1$ 为 $\lambda J_1(\lambda)=Bi_m J_0(\lambda)$ 第一根（圆柱 Robin
特征值问题，标准谱理论见~\cite{crank}），$A_1$ 为中心响应系数。
问题 3（式~\eqref{eq:时标}）：$D_{\max}=1.359\times10^{-8}$ m$^2$/s、$Bi_m=1.177$，得
$t_{\mathrm{ref}}=15.70$ h（300 模精确级数穿越同为 15.7021 h，参考问题自身的
首模近似误差可忽略）；膜界 $11.17$ h（忽略内部扩散阻力，系统性偏弱，
仅作次级参照）。问题 4：沿真实 $R(s)$ 轨迹的瞬时首模时间积分版参考为
$19.02$ h（$R_{\min}$ 常值版 15.95 h、固定 $R_0$ 版 37.61 h）。
两个主值（57.18 h、50.54 h）均远大于各自参考时标。

\textbf{这些数字是参考时标，不是严格误差界}：(i) 参考解把环境平衡势固定为
末值 $C_e=0.04986$，但真实边界初期低至 0.01963——同一 $D_{\max}$ 首模参考
在固定初值边界下为 \textbf{14.5491} h，"排除线"随边界约定移动，比较序不自动成立；
(ii) 对散度形式的变系数扩散方程，$D\le D_{\max}$ 本身不构成逐点解比较；
(iii) 收缩版积分瞬时第一特征值，而 $Bi(s)$ 随时间变化引起模态耦合，
瞬时特征值积分不是精确的变边界谱解。远低于参考时标的达标时间应触发怀疑，
但它们不构成已证明的严格下界。

\subsection{六层验证体系}
可靠性来自\textbf{两种不同数学结构的互证}，而不是同一算法换网格。
六层验证体系见表~\ref{tab:验证}（V1--V6 为本文自定义分层编号），
每层写明能排除什么、不能排除什么；收敛与守恒的可视化见图~\ref{fig:收敛守恒}。
两点纪律：(i) V1 与 V5 的 $D$ 冻结档是解析结构对有限体积解的\emph{真正独立}比较
（不是换参数、不是换网格；剖面对比见图~\ref{fig:对拍}）；V2 为后验尾段诊断，
不列入独立验证；(ii) V4/V6 属于一致性检查，只能证明实现与离散无误，
不把它们说成可靠性证明。

\begin{table}[H]
\centering
\caption{验证汇总表（V1--V6）：手段、数学结构、实测数字、可排除与不能排除的错误。}
\label{tab:验证}
\footnotesize
\begin{tabular}{p{0.30\linewidth}p{0.30\linewidth}p{0.32\linewidth}}
\toprule
层（手段） & 实测数字 & 能排除 / 不能排除 \\
\midrule
V1 问题 1 热场解析谱解（特征展开 + Duhamel 卷积）vs 有限体积数值解 &
单元中心逐点最大偏差 $5.9\times10^{-7}\,^{\circ}$C（N=320）；收敛阶 2.00/1.99/1.99 &
排除离散、边界处理、时间积分错误（两种完全不同方法）；不能排除材料参数取值错误 \\
V2 后验尾段诊断（等温线性化 Bessel 首模结构对照） &
尾段 $\ln(U_{\max}-C_e)$ 线性 $R^2=0.9972$；衰减率 $\gamma=3.55\times10^{-6}$/s 对应
等效 $D=2.54\times10^{-10}$（$C\approx0.108$，在尾段剖面区间内） &
尾段结构诊断与等效参数解释（$D_{\mathrm{eff}}$ 由待验证解反演，属构造一致，
不是独立预测）；不能排除一般的边界或几何离散错误 \\
V3 参考时标门槛检验（常系数/瞬时首模近似，见 6.2 节限定） &
$t_f=57.18$ h $>$ 15.70 h；$t_f=50.54$ h $>$ 19.02 h &
远低于参考时标的达标时间应触发怀疑；非严格误差界，不能给出上界 \\
V4 守恒恒等式（总水量、冻结热容的显热离散平衡，逐步对账） &
水量窗口相对残差 $\le5.5\times10^{-11}$；显热离散平衡全局吞吐残差 $\le1.1\times10^{-11}$ &
核对离散线性系统的求解一致性；未独立定义状态焓，
不证明真实能量闭合，也不能替代潜热分析 \\
V5 退化测试（移动域 $R(t)\equiv R_0$；$D$ 冻结） &
移动域退化为固定域：$\Delta t_f=0.00$ s、场差 0；$D$ 冻结退回 Bessel 首模：偏差
$3.1\times10^{-9}$、首模商 0.99996 &
排除移动域实现错误；常系数湿场退化为真独立对照；不能排除闭合假设本身 \\
V6 网格/时间 Richardson 收敛 &
场级空间阶 1.90--2.45；$t_f$ 序列 57.249/57.201/57.180 h，外推偏差 0.016 h；
rtol/10 实测 $\Delta t_f=0.0014$ h &
只证明离散收敛（一致性）；不能单独充当可靠性证据 \\
\bottomrule
\end{tabular}
\end{table}

\begin{figure}[htbp]
\centering
\includefig[0.86\linewidth]{fig2_问题1解析对拍.pdf}{问题 1 解析对拍}
\caption{问题 1 解析比较（V1）：(a) 特征位置（中心/表面）温度的
数值解与解析谱解时程；(b) 两个特征位置的逐点偏差时程，量级
$10^{-6}\,^{\circ}$C、全程低于 $3.2\times10^{-6}\,^{\circ}$C，虚线为四位小数输出的
分辨率上限；径向网格加密的二阶收敛见表~\ref{tab:验证} 与
图~\ref{fig:收敛守恒}(a)；全域单元中心在 N=320 的最大偏差为
$5.9\times10^{-7}\,^{\circ}$C，即正文引用的 V1 验证量。}
\label{fig:对拍}
\end{figure}

\begin{figure}[htbp]
\centering
\includefig[\linewidth]{fig15_收敛与守恒检验.pdf}{收敛与守恒检验}
\caption{收敛与守恒检验。(a) 收敛·温度场误差：解析比较偏差
随网格的二阶下降（V1/V6）；(b) 收敛·$t_f$ 随网格的收敛序列；
(c) 守恒·残差随网格加密：水量与显热离散平衡残差维持在 $10^{-11}$ 量级附近（V4）。}
\label{fig:收敛守恒}
\end{figure}

\section{结果分析与敏感性}

交付表格（表~\ref{tab:表1}--表~\ref{tab:表6}，对应 result1--4.xlsx）汇总如下；
达标穿越的三线对照（问题 3 主解、问题 4 主解、附录 4 固定半径对照）见图~\ref{fig:穿越}。

% 本文件由 gen_tables.py 从 03-数据/*.csv 自动生成，禁止手工改数。

\begin{table}[htbp]
\centering
\caption{问题~1（预热平衡阶段）药材温度（℃）。数字由 result 工作簿同一数据源抽取，保留四位小数。}
\label{tab:表1}
\small
\begin{tabular}{lccccc}
\toprule
时间/s & 0 & 0.5 & 1 & 1.5 & 2 \\
\midrule
100s & 28.0001 & 28.0003 & 28.0040 & 28.0327 & 28.1801 \\
300s & 28.0408 & 28.0635 & 28.1514 & 28.3681 & 28.8489 \\
600s & 28.4533 & 28.5360 & 28.8039 & 29.3159 & 30.1652 \\
900s & 29.3243 & 29.4583 & 29.8755 & 30.6161 & 31.7304 \\
1200s & 30.5427 & 30.7098 & 31.2223 & 32.1126 & 33.4276 \\
1500s & 31.9957 & 32.1867 & 32.7660 & 33.7463 & 35.1203 \\
1800s & 33.5753 & 33.7720 & 34.3642 & 35.3621 & 36.7856 \\
\bottomrule
\end{tabular}
\end{table}

\begin{table}[htbp]
\centering
\caption{问题~1（预热平衡阶段）药材水分浓度（kg/kg）。数字由 result 工作簿同一数据源抽取，保留四位小数。}
\label{tab:表2}
\small
\begin{tabular}{lccccc}
\toprule
时间/s & 0 & 0.5 & 1 & 1.5 & 2 \\
\midrule
100s & 2.5500 & 2.5500 & 2.5500 & 2.5500 & 2.2470 \\
300s & 2.5500 & 2.5500 & 2.5500 & 2.5492 & 2.0508 \\
600s & 2.5500 & 2.5500 & 2.5500 & 2.5353 & 1.8770 \\
900s & 2.5500 & 2.5500 & 2.5497 & 2.5045 & 1.7547 \\
1200s & 2.5500 & 2.5500 & 2.5482 & 2.4646 & 1.6586 \\
1500s & 2.5500 & 2.5499 & 2.5445 & 2.4206 & 1.5787 \\
1800s & 2.5500 & 2.5497 & 2.5383 & 2.3755 & 1.5102 \\
\bottomrule
\end{tabular}
\end{table}

\begin{table}[htbp]
\centering
\caption{问题~2（全烘干过程）药材温度（℃）。数字由 result 工作簿同一数据源抽取，保留四位小数。}
\label{tab:表3}
\small
\begin{tabular}{lccccc}
\toprule
时间/h & 0 & 0.5 & 1 & 1.5 & 2 \\
\midrule
0.5h & 32.1892 & 32.3821 & 32.9659 & 33.9608 & 35.4130 \\
1.0h & 40.3816 & 40.5536 & 41.0605 & 41.8782 & 42.9976 \\
1.5h & 45.8468 & 45.9348 & 46.1930 & 46.6051 & 47.1401 \\
2.0h & 48.4502 & 48.4881 & 48.5984 & 48.7735 & 49.0033 \\
2.5h & 49.4670 & 49.4792 & 49.5136 & 49.5653 & 49.6609 \\
3.0h & 49.8495 & 49.8553 & 49.8746 & 49.9101 & 49.9664 \\
\bottomrule
\end{tabular}
\end{table}

\begin{table}[htbp]
\centering
\caption{问题~2（全烘干过程）药材水分浓度（kg/kg）。数字由 result 工作簿同一数据源抽取，保留四位小数。}
\label{tab:表4}
\small
\begin{tabular}{lccccc}
\toprule
时间/h & 0 & 0.5 & 1 & 1.5 & 2 \\
\midrule
0.5h & 2.5499 & 2.5489 & 2.5256 & 2.3257 & 1.6486 \\
1.0h & 2.5258 & 2.4948 & 2.3578 & 2.0230 & 1.4711 \\
1.5h & 2.3861 & 2.3256 & 2.1344 & 1.8020 & 1.3476 \\
2.0h & 2.1709 & 2.1084 & 1.9236 & 1.6259 & 1.2311 \\
2.5h & 1.9566 & 1.9006 & 1.7360 & 1.4720 & 1.1166 \\
3.0h & 1.7662 & 1.7165 & 1.5702 & 1.3333 & 1.0081 \\
\bottomrule
\end{tabular}
\end{table}

\begin{table}[htbp]
\centering
\caption{问题~3 药材烘干过程的水分浓度（kg/kg）。结束行时刻为未舍入判定值 57.1799~h；表中 0.1500 为舍入显示，判定位用未舍入值。}
\label{tab:表5}
\small
\begin{tabular}{lccccc}
\toprule
时间/h & 0 & 0.5 & 1 & 1.5 & 2 \\
\midrule
6h & 1.0156 & 0.9868 & 0.9006 & 0.7546 & 0.5335 \\
12h & 0.4548 & 0.4418 & 0.4019 & 0.3289 & 0.1639 \\
18h & 0.2979 & 0.2906 & 0.2679 & 0.2247 & 0.0842 \\
24h & 0.2371 & 0.2321 & 0.2161 & 0.1851 & 0.0663 \\
30h & 0.2053 & 0.2013 & 0.1887 & 0.1637 & 0.0597 \\
36h & 0.1854 & 0.1821 & 0.1714 & 0.1501 & 0.0565 \\
42h & 0.1717 & 0.1687 & 0.1594 & 0.1404 & 0.0547 \\
48h & 0.1615 & 0.1588 & 0.1504 & 0.1332 & 0.0536 \\
54h & 0.1535 & 0.1511 & 0.1434 & 0.1275 & 0.0528 \\
烘干结束时间 & 0.1500 & 0.1477 & 0.1402 & 0.1249 & 0.0525 \\
\bottomrule
\end{tabular}
\end{table}

\begin{table}[htbp]
\centering
\caption{问题~4 药材烘干过程的水分浓度（kg/kg）。“—”表示该时刻该固定位置已在药材表面之外（半径收缩，留空不外推）。结束行时刻为未舍入判定值 50.8230~h。}
\label{tab:表6}
\small
\begin{tabular}{lccccc}
\toprule
时间/h & 0 & 0.5 & 1 & 1.5 & 药材表面 \\
\midrule
6h & 1.7164 & 1.5350 & 1.0212 & — & 0.4215 \\
12h & 0.7340 & 0.6516 & 0.4068 & — & 0.1666 \\
18h & 0.4065 & 0.3667 & 0.2387 & — & 0.0888 \\
24h & 0.2837 & 0.2598 & 0.1783 & — & 0.0671 \\
30h & 0.2253 & 0.2085 & 0.1488 & — & 0.0593 \\
36h & 0.1920 & 0.1790 & 0.1312 & — & 0.0558 \\
42h & 0.1706 & 0.1598 & 0.1196 & — & 0.0539 \\
48h & 0.1556 & 0.1463 & 0.1113 & — & 0.0529 \\
烘干结束时间 & 0.1500 & 0.1413 & 0.1081 & — & 0.0525 \\
\bottomrule
\end{tabular}
\end{table}

\begin{figure}[htbp]
\centering
\includefig[0.86\linewidth]{fig4_达标穿越.pdf}{达标穿越}
\caption{达标穿越：$\max_qU(q,t)$ 全程曲线与
0.15 阈值线；问题 3 主解（$t_f=57.1799$ h）、问题 4 主解（$t_f=50.5369$ h，
内生几何）与附录 4 固定半径对照（129.1047 h）三线对比，横轴标出参考时标
（15.70 h / 19.02 h，常系数/瞬时首模近似，非严格误差界）。}
\label{fig:穿越}
\end{figure}

\subsection{敏感性分析}
每组都实测 $\Delta t_f$（不只报方向），见表~\ref{tab:敏感性} 与图~\ref{fig:敏感}。
边界汇取值（$C_e$）是本题\emph{唯一能改变结论量级的单点}：改用解吸等温线
$C_e=0.1366$ 后 $t_f$ 增大 $22.6$ h（$39.6\%$）。按既定做法，主结果一律用声明的
闭合设定，题外修正只进入敏感性分析并明确标注；该偏移不计入误差预算（7.2 节）。
半径插值两档实测仅差 $0.0022$ h（附件 2 平滑，两插值几乎重合），说明该设定不敏感。

\begin{table}[htbp]
\centering
\caption{敏感性矩阵（问题 3 设定、N=80 对照档；问题 4 拆分用主档）。}
\label{tab:敏感性}
\small
\begin{tabular}{lll}
\toprule
组别 & $t_f$ (h) & $\Delta t_f$ (h) \\
\midrule
$C_e=0.04986$（声明闭合，主） & 57.2009 & — \\
$C_e=0.1366$（山药解吸支，题外修正） & 79.8270 & $+22.6261$（$+39.6\%$） \\
环境外推：末值保持（主） & 57.2009 & — \\
环境外推：线性外推（末 1 h 斜率） & 58.8203 & $+1.6194$ \\
环境外推：末 1 h 均值 & 57.5034 & $+0.3025$ \\
半径插值：PCHIP（主，问题 4，附件 2 外生几何对照） & 50.8295 & — \\
半径插值：线性（问题 4） & 50.8317 & $+0.0022$ \\
问题 4 拆分：附录 4 固定 $R$ & 129.1047 & $+78.57$（对内生主解） \\
问题 4 拆分：附录 4 + 内生 $R(t)$（主解） & 50.5369 & — \\
问题 4 拆分：问题 3 & 57.1799 & $+6.64$（对内生主解） \\
\bottomrule
\end{tabular}
\end{table}

\begin{figure}[htbp]
\centering
\includefig[\linewidth]{fig5_敏感性.pdf}{敏感性}
\caption{敏感性条形图：各因素 $\Delta t_f$ 横向条形，
边界汇取值 $C_e=0.1366$（题外修正）$+22.63$ h 自然主导，环境外推 $+1.62$ h 次之，
半径插值仅 $+0.002$ h；图中给出最大不确定性来源，主结果一律用声明的闭合设定。}
\label{fig:敏感}
\end{figure}

\subsection{误差与不确定度（三类分别报告）}
误差预算按\textbf{数值误差、环境情景、模型结构差异}三类分别报告
（表~\ref{tab:预算}），问题 3/4 分列；其含义统一为
"在所列模型假设与情景集合内，已检查因素的累计偏移量级"，
不夸大为实际烘干时长的可信区间，也不是已证明的严格误差界。

\begin{table}[htbp]
\centering
\caption{误差预算表（三类分别报告、两问分列；$C_e$ 取值与潜热情景属模型结构差异，
单列不合成）。}
\label{tab:预算}
\small
\begin{tabular}{lp{0.56\linewidth}ll}
\toprule
类别 & 误差源 & 问题 3 (h) & 问题 4 (h) \\
\midrule
数值误差 & 空间网格（Q3：N=160 Richardson 偏差；Q4：连续事件定位 N=80→160 变化） & 0.0164 & 0.0007 \\
数值误差 & 时间步（rtol/10 实测） & 0.0014 & 0.0014 \\
数值误差 & 非线性迭代（Picard/Newton 对照） & $1\times10^{-6}$ & $1\times10^{-6}$ \\
数值误差 & 输出重建 / 1 s 插值（曲率上界实测） & $2.6\times10^{-4}$ & $2.6\times10^{-4}$ \\
环境情景 & 环境外推（三档最大偏移） & 1.6194 & 1.6194 \\
设定依赖 & 半径插值（线性 vs PCHIP，仅问题 4） & — & 0.0022 \\
设定依赖 & 舍入判定（60 s 网格 + 四位小数，另一种数据使用方式） & 0.0701 & — \\
\midrule
\multicolumn{2}{l}{\textbf{已检查因素累计偏移量级（三角和）}} & $\approx1.71$ & $\approx1.62$ \\
\bottomrule
\end{tabular}
\end{table}

\textbf{模型结构差异（单列，不计入上表合成）}：潜热闭合情景 $+5.3889$ h
（5.3 节，情景假设集合下无潜热基线的 $t_f$ 偏差量级）；
边界汇取值 $C_e=0.1366$ 时为 $+22.6261$ h（题外修正，单列敏感性）。
\textbf{未覆盖项清单}：边界闭合的全部允许形式（吸附/水活度关系缺数据）、
内生收缩标定参数（$k_p$ 的 95\% CI 相对半宽 24\%）、内部非仿射变形
（5.5 节的结构性分析，不计入主线预算）。
故问题 3 的达标时间报告为 $t_f=57.18$ h，在所列模型假设与情景集合内，
已检查因素的累计偏移量级约 1.71 h（其中数值误差仅占 $0.02$ h 量级，
主导项是环境外推的取值方式）。

\section{模型评价与改进}

\subsection{优点}
\begin{enumerate}
  \item \textbf{统一性与自洽}：四问为一套热湿耦合方程的四个配置（表~\ref{tab:开关}），
        问题 1 是问题 2 的解耦特例、问题 4 是移动域推广，无模型拼接缝隙；
        质量坐标下收缩几何方程同形、无伪对流项（附录 A 双守恒推导）。
  \item \textbf{验证分层且独立}：解析谱解与守恒有限体积两种不同数学结构互证
        （V1 逐点 $5.9\times10^{-7}\,^{\circ}$C、二阶收敛），叠加退化测试、
        守恒恒等式（$10^{-11}$ 量级）与参考时标门槛检验，每层都明确能够排除
        与不能排除的错误类型。
  \item \textbf{收缩由模型内生预测}：问题 4 不读附件 2 而由半经验闭合逐步预测
        $R(t)$（$t_f=50.5369$ h），附件 2 降为验证数据后仍全程吻合
        （0.088 mm），两个精度指标与标定循环的局限均已定量说明。
  \item \textbf{不确定度透明}：判据族反验、敏感性全 $\Delta t_f$、误差三类分别报告、
        潜热与边界闭合两个情景各自单列，结论的适用边界可读可查。
\end{enumerate}

\subsection{缺点}
\begin{enumerate}
  \item \textbf{湿边界是声明的有效闭合}：$C_{\mathrm{env}}$ 直接作 Robin 驱动势
        缺吸附等温线支撑，是本题最大单点依赖（题外修正 $+22.6261$ h）。
  \item \textbf{无潜热基线}：潜热与显热同量级，基线的 $t_f$ 偏差量级约
        $5.4$ h（表面相变汇情景），更真实的闭合需要题目未给的吸附/解吸数据。
  \item \textbf{收缩闭合是半经验}：$g_c,g_p$ 为经验形状函数，力学驱动扩展虽
        更严谨，但在所选幂律保水—简化储能结构内存在精度上限
        （末端 $J=0.3343$ vs 0.36），缺保水曲线与力学实测数据。
  \item \textbf{几何与外推的理想化}：一维径向（端面效应未给定量界）、
        环境 4 h 后末值保持、内部等比收缩（内部收缩不可识别）。
\end{enumerate}

\subsection{改进}
\begin{enumerate}
  \item 针对缺点 1/2：补测材料吸附/解吸等温线（或水活度关系）与相变焓数据，
        把湿边界与能量边界升级为介质配套的闭合；
  \item 针对缺点 3：补测真实保水曲线 $p_c(S,T)$ 与骨架模量/松弛谱，
        沿力学驱动路线（5.5 节骨架已经建立）重新标定，半经验闭合作为低维验证基准；
  \item 针对缺点 4：补有限长圆柱对照量化端面效应、以更长时程环境数据替换外推、
        在力学框架内放松仿射假设（非仿射已在扩展模型中实现）；
  \item 若保留严格数学亮点：固定全部比较条件后构造可证明的包络或积分不等式，
        把参考时标升级为已证明的严格下界。
\end{enumerate}

\section*{参考文献}
\begingroup
\renewcommand{\section}[2]{}% 抑制 thebibliography 自动 \refname，避免双标题
\begin{thebibliography}{9}
\bibitem{luikov} Luikov A V. \emph{Heat and Mass Transfer in Capillary-Porous Bodies}.
Pergamon Press, Oxford, 1966.
\bibitem{crank} Crank J. \emph{The Mathematics of Diffusion}. 2nd ed.,
Oxford University Press, 1975.
\bibitem{pharmacopoeia} 国家药典委员会. \emph{中华人民共和国药典（一部）}.
北京: 中国医药科技出版社, 2020.（药材水分限度的通行标准约 13\%；本文未核对具体条文与版本卷次）
\bibitem{saravacos} Saravacos G D. Mass transfer properties of foods. In:
\emph{Engineering Properties of Foods} (M. A. Rao, S. S. Rizvi, eds.),
Marcel Dekker, New York, 1986.
\end{thebibliography}
\endgroup

\appendix
\ctexset{section = {name = {附录 ,：}, number = \Alph{section},
           format = \centering\heiti\zihao{4}\bfseries}}

\section{移动域方程的推导要点}
正确的出发点必须同时包含干物质守恒与水分守恒。设 $\rho_d$ 为单位总体积中的
干物质质量、$v$ 为骨架径向速度、$j$ 为相对骨架的水通量（$j=-\rho_d D\,\partial_r U$）：
\begin{equation}
\partial_t\rho_d+\frac1r\partial_r(r\rho_d v)=0,\qquad
\partial_t(\rho_d U)+\frac1r\partial_r\!\left[r(\rho_d Uv+j)\right]=0 .
\end{equation}
第二式展开后用第一式相减，得物质点形式
\begin{equation}
\rho_d\,(U_t+vU_r)=\frac1r\partial_r\!\left(r\rho_d D\,U_r\right).
\label{eq:物质点}
\end{equation}
引入材料坐标 $q$（干物质标签），物理坐标 $r=\psi(q,t)$，骨架速度
$v=\psi_t|_q$。在\textbf{干物质密度空间均匀}与\textbf{仿射收缩}
$\psi=R(t)\sqrt q$（长度不变、内部等比；附件 2 只识别外半径，
内部收缩不可识别，等比是最简闭合）假设下，$\rho_d\psi\psi_q$ 与 $t$ 无关，
式~\eqref{eq:物质点} 化为材料坐标形式
$U_t|_q=\partial_q\!\left(4qD/R(t)^2\,U_q\right)$——与固定域同形、无伪对流项。
在移动表面 $q=1$ 处，对表面薄层做水分衡算，几何运动项与通量跳条件中的对应项
恒等相消，Robin 边界 $-D\partial_rU=h_m(U_s-C_{\mathrm{env}})$ 形式不变。温度方程同理。
两点说明：(i) 旧稿"从 $U_t|_r=(1/r)\partial_r(rDU_r)$ 出发、迁移项相消"
的写法不成立（该式右侧没有可相消项），以上双守恒推导才是正确路径；
(ii) 固定半径退化测试（V5a）\textbf{无法检验}这一缺口：
此时骨架速度 $v\equiv0$，移动域与固定域方程逐字相同，
它只能检验实现的一致性，不能检验移动域推导本身。

\section{参数取值表与收缩参数表}
全部非常数参数\textbf{在哪个位置、哪个时刻、用哪个公式、什么单位}取值，
见表~\ref{tab:参数}；内生收缩模型的参数（拟合/非拟合分栏）见表~\ref{tab:收缩参数}；
符号的量纲逐条核对见表~\ref{tab:量纲}。

\begin{table}[H]
\centering
\caption{参数取值表（题给 + 声明假设分栏）。}
\label{tab:参数}
\small
\begin{tabular}{llll}
\toprule
参数 & 取值/公式 & 取值方式 & 适用 \\
\midrule
$D(C)$ & $7\times10^{-9}\mathrm e^{-0.89/C}$ & 局部 $C$ & 问题 1 \\
$D(C,T)$ & $2.4\times10^{-3}\mathrm e^{-0.45/C}\mathrm e^{-3850/T}$ & 局部 $C$、局部 $T$（开尔文） & 问题 2/3 \\
$D(C,T)$ & $4.2\times10^{-4}\mathrm e^{-0.30/C}\mathrm e^{-3850/T}$ & 局部 $C$、局部 $T$（开尔文） & 问题 4 \\
$h$ & $25$ W/(m$^2\cdot$K) & 常数 & 问题 1--4 \\
$h_m$ & $8\times10^{-7}$ m/s & 常数 & 问题 1--4 \\
$T_a(t),C_{\mathrm{env}}(t)$ & 附件 1（0--14400 s，60 s 步） & 分段线性插值 & 问题 1--4 \\
$R(t)$ & 附件 2（0--259200 s，1800 s 步） & PCHIP 保形插值 & 问题 4 \\
几何 & $L_0=0.25$ m，$R_0=0.02$ m & 一维径向（$L_0/R_0=12.5$） & 问题 1--4 \\
初值 & $U_0=2.55$ kg/kg，$\varTheta_0=301.15$ K & 全域均匀 & 问题 1--4 \\
\midrule
\multicolumn{4}{l}{\textit{声明假设（题目未给定，本文声明并前置）}} \\
湿边界闭合 & $C_{\mathrm{env}}$ 直接作 Robin 驱动势 & 有效闭合（第三章假设 2） & 问题 1--4 \\
环境外推 & $t>14400$ s 末值保持 & 声明（敏感性见 7.1 节） & 问题 2--4 \\
$h,h_m$ 延用 & 附录 2 值延用到问题 2--4 & 声明为参数延用假设 & 问题 2--4 \\
半径延拓 & $t>259200$ s 保持 1.198 cm & 声明 & 问题 4 \\
\bottomrule
\end{tabular}
\end{table}

\begin{table}[H]
\centering
\caption{内生收缩模型参数表（拟合/非拟合分栏标注）与物理合理性。}
\label{tab:收缩参数}
\small
\begin{tabular}{p{0.22\linewidth}lp{0.28\linewidth}p{0.24\linewidth}}
\toprule
参数 & 数值 & 来源 & 物理合理性核对 \\
\midrule
$\varepsilon_c$（塌陷峰值应变） & 0.1108 & 拟合附件 2 & 11.1\% 在毛细压实区间 $[1.4\%,29\%]$ 内 \\
$a,\ b$（塌陷形状指数） & 1.668,\ 1.220 & 拟合附件 2 & 峰值相位 $\bar U\approx1.47$（约 4 h）与急缩段一致 \\
$\varepsilon_p$（停滞幅度） & 0.0709 & 拟合附件 2 & 对应末期孔隙率 30.9\% $\approx$ 附件 2 隐含 31.0\% \\
$k_p$（停滞开启陡峭度） & 3.787 & 拟合附件 2 & 开启相位由 $C_s$ 穿越 $C_{\mathrm{glass}}$ 物理锚定 \\
$C_{\mathrm{glass}}$（玻璃化含水率） & 0.21 & Gordon--Taylor 结构典型参数情景值，未拟合 & 估计区间 0.19--0.23，列入参数不确定清单 \\
$\rho_{sk}$（骨架密度，kg/m$^3$） & 1468 & 生物质骨架密度典型范围值，未拟合 & 结果对其不敏感 \\
\bottomrule
\end{tabular}
\end{table}

\begin{table}[H]
\centering
\caption{量纲检查表（全部符号逐条核对）。}
\label{tab:量纲}
\small
\begin{tabular}{lll}
\toprule
符号 & 单位 & 核对要点 \\
\midrule
$U,\ C$ & kg/kg（干基） & 边界驱动势 $C_{\mathrm{env}}$ 同基准（声明闭合） \\
$D$ & m$^2$/s & $4qD/R^2$ 为 1/s，与 $U_t$ 一致 \\
$T$ & K（不是 ${}^\circ$C） & $\mathrm e^{-3850/T}$ 的 3850 单位为 K；读入与全程开尔文防呆断言 \\
$h,h_m$ & W/(m$^2\cdot$K)、m/s & $-k\partial_rT=h(T_s-T_a)$ 两侧 W/m$^2$；$h_m\Delta U$ 为 m/s \\
$a_{\mathrm{eff}}$ & J/(m$^3\cdot$K) & $\rho c_p$ 逐点取值，$\rho$ 不进湿分方程 \\
$q$ & — & $q=(X/R_0)^2\in[0,1]$，干物质标签 \\
$\varepsilon_c,\varepsilon_p$ & —（应变） & 式~\eqref{eq:收缩闭合} 的乘子修正 \\
\bottomrule
\end{tabular}
\end{table}

\section{内生收缩模型的标定细节}
内生收缩模型（式~\eqref{eq:收缩闭合}）的 5 个参数
（$\varepsilon_c,a,b,\varepsilon_p,k_p$）由 scipy 的 \texttt{least\_squares}
在附件 2 的 145 个采样点上对 $R_{\mathrm{pred}}-R_{\mathrm{data}}$ 做带物理边界约束的
最小二乘估计；$C_{\mathrm{glass}}=0.21$（Gordon--Taylor 结构典型参数情景值）与
$\rho_{sk}=1468$ kg/m$^3$（生物质骨架密度典型范围值）不参与拟合，均列入参数不确定清单。
过拟合自检用奇偶折半留出法：两折上训练与留出 RMSE 均在 0.045--0.050 mm，
且两折参数互差 $<10\%$——但该检验只覆盖同一光滑曲线上的插值
（湿度场由附件 2 实测几何生成，存在标定循环），只能说明插值稳健，
不构成参数被独立识别的证明；块 bootstrap（残差自相关块长 24、$B=200$）95\%
置信区间的相对半宽 $\varepsilon_c$ 4.7\%、$a$ 11.9\%、$b$ 11.0\%、
$\varepsilon_p$ 3.8\%、$k_p$ 23.9\%（分布有偏斜/多峰迹象）；
变初值 $\times$ 变损失函数的最大参数偏差 $\le0.98\%$；
各机制在其主导区段单独拟合与全曲线拟合的幅度参数差 $<2\%$，真正的独立识别需完整耦合标定。
\textbf{参数的第二路线检验}：$\varepsilon_c=11.1\%$ 落在毛细压
$2\gamma/r_{\mathrm{pore}}$ 对基体模量的理论应变区间 $[1.4\%,28.8\%]$；
$\varepsilon_p$ 对应的末期孔隙率 30.9\% 与"平台半径+质量守恒"独立推出的
31.0\% 一致；$C_{\mathrm{glass}}=0.21$ 与"附件 2 冻结起点约 21 h
$\leftrightarrow$ 模型 $\bar U(21\ \mathrm{h})=0.234$"交叉吻合；
$\rho_{sk}=1468$ 在生物质骨架密度典型范围内且结果对其不敏感。
第二路线检验的结果是：$\varepsilon_c$、$a$、$b$、$\varepsilon_p$、$C_{\mathrm{glass}}$、
$\rho_{sk}$ 均获得支持；\textbf{$k_p$ 存疑}——只有曲线拟合一条路
（无独立物理观测量、置信区间最宽 24\%），但它只影响过渡带形状，
不触及达标时间与平台半径的结论。
力学扩展模型的标定（4 起点同一解、留出、bootstrap、末端平衡分析）
见 5.5 节正文与数据文件 \texttt{03-数据/mechRefit.csv}。

\section{复现清单与代码结构}
一条命令 \texttt{python run\_all.py}（Windows 下 \texttt{PYTHONIOENCODING=utf-8}，
依赖仅 numpy/scipy/pandas/openpyxl）依次执行：
\begin{itemize}
  \item \texttt{run\_q1.py}：问题 1 求解 + V1 解析比较 + 守恒与退化检验 + result1.xlsx；
  \item \texttt{run\_q23.py}：问题 2/3 耦合求解 + $t_f$ + V3/V4/V6 门槛检验 +
        result2/result3.xlsx + 环境外推三档；
  \item \texttt{run\_q4.py}：问题 4 移动域 + 效应拆分 + V5 退化 + result4.xlsx；
        \texttt{run\_q4\_endo.py}：内生几何主解（满 72 h 积分、三段误差评价）；
  \item \texttt{run\_m4.py}：判据族反验 + 敏感性矩阵 + 误差预算 + V2 尾段 + 结果总账；
  \item \texttt{latent\_scenario.py}：潜热对照情景（5.3 节）；
  \item \texttt{run\_mech\_refit.py}（03-数据）：力学扩展模型重标定（5.5 节）。
\end{itemize}
全部交付文件确定性写出（固定归档时间戳），连续两次运行 result1--4 的 md5
完全一致（result1 \texttt{22eb47f7…}、result2 \texttt{12cde4e8…}、
result3 \texttt{fbe982ff…}、result4 \texttt{80c73431…}）。
核心求解器源码摘录（守恒步进与表面重建，注释从略）如下。

\begin{lstlisting}[language=Python,caption={守恒耦合步进核心（solver 摘录）}]
# 表面值：半单元内阻与外对流阻串联（D_eff 取半单元中点值）
def surface_m(CN, Ts, Cenv, dq):
    g = lambda Cs: 8.0*D_of(0.5*(Cs+CN), Ts)*(Cs-CN)/(R0**2*dq) \
                   - 2.0*HM*(Cenv-Cs)/R0
    Cs = brentq(g, lo, hi, xtol=1e-15, rtol=1e-14)
    return Cs, 2.0*HM*(Cenv-Cs)/R0

# 一个耦合后向欧拉步：热场（冻结系数，线性）→ 湿场（冻结温度，Newton）
def coupled_step(U, T, dt, Ta, Cenv, dq, newton):
    ch = 4.0*qf*k_of(Uf)*dt/(R0**2*dq**2)      # 面通量系数
    Rh = R0**2*dq/(8.0*ks) + R0/(2.0*H)        # 表面串联热阻
    T_new = solve_banded((1,1), ab_T, rhs_T)   # 热场线性求解
    for _ in range(30):                        # 湿场 Newton/Picard
        dU = solve_banded((1,1), J, -G)
        U_new = U_new + dU
    return U_new, T_new
\end{lstlisting}
\end{document}
